\section{Discussion}
\label{sec:discussion}

The results presented in this technical report demonstrate the effectiveness of MAC in enabling coordinated problem-solving through multiple specialized agents. Across the evaluated domains of Travel, Mortgage, and Software Development, the multi-agent collaboration approach achieved impressive goal success rates, with the overall GSR reaching as high as 90\% when utilizing the Claude 3.5 Sonnet (20241022) models for both the supervisor agent and specialist agents.
This strong performance is particularly notable when compared to the single-agent setting, where a significant regression of up to 37\% was observed in the Software Development domain. The ability of the multi-agent collaboration to leverage the specialized expertise of individual agents appears to be a key factor in this improved effectiveness.

Further analysis of the results highlights the importance of optimizing the inter-agent communication mechanisms. The impact of the payload referencing feature is most pronounced in the Software Development domain, where it led to a 23\% relative improvement in overall GSR as well as a 27\% relative reduction in the average communication overhead per turn. By allowing the supervisor agent to more efficiently reference and share large content blocks, such as code snippets, this specialized mechanism enhances the reliability and speed of the multi-agent interactions.
However, the results also reveal some limitations of the current system. The higher latency observed in the more complex Software Development scenarios suggests that further optimizations may be needed to reduce the overhead of multi-agent coordination, particularly in time-sensitive applications. 
